\documentclass[10pt,a4paper]{article}
\usepackage[utf8]{inputenc}
\usepackage{amsmath}
\usepackage{amsfonts}
\usepackage{amssymb}
\usepackage{float}
\usepackage{enumerate}
\title{GSF-3100 \\
04 Structure des taux d'intérêt \\
Série d'exercices 5}

\begin{document}
\maketitle

\begin{itemize}
\item Dans ce document vous avez des questions à choix multiples allant de a) à d).  
\item Il y a seulement un choix possible par question.
\item Ce questionnaire contient des exercices en lien avec la structure des taux d'intérêt.
\item Contient 17 questions. 
\end{itemize}
 
\newpage

\section*{Section 1}
Considérer les informations suivantes sur des obligations gouvernementales:


\begin{table}[H]
\begin{center}
\begin{tabular}{|l|l|l|l|}
\hline
Période t & Échéance (an) & TC (\%, m = 2) & y (\%, m = 2) \\ \hline
1         & 0,5           & 0              & 6,00          \\ \hline
2         & 1             & 0              & 5,80          \\ \hline
3         & 1,5           & 5,75           & 5,55          \\ \hline
4         & 2             & 6,00           & 5,25          \\ \hline
\end{tabular}
\end{center}
\end{table}



\subsection*{Exercice 1}
Pour l'obligation d'échéance 1,5 an (18 mois),  trouver la prime de risque (ou l’écart à terme).


\begin{enumerate}[a)]
\item -0.40\%
\item -0.45\%
\item -0.50\%
\item -0.55\%
\end{enumerate}

\subsection*{Exercice 2}
Pour l'obligation d'échéance 1,5 an (18 mois),  trouver la prime de risque relative.


\begin{enumerate}[a)]
\item -6.00\%
\item -6.50\%
\item -7.00\%
\item -7.50\%
\end{enumerate}

\subsection*{Exercice 3}
Pour l'obligation d'échéance 1,5 an (18 mois),  trouver le ratio de taux (en nombre de fois)


\begin{enumerate}[a)]
\item 1.025
\item 1.00
\item 0.925
\item 0.900
\end{enumerate}

\section*{Section 2}
Considérer les informations suivantes sur des obligations gouvernementales:

\begin{table}[H]
\begin{center}
\begin{tabular}{|l|l|l|l|}
\hline
Période t & Échéance (an) & TC (\%, m = 2) & y (\%, m = 2) \\ \hline
1         & 0,5           & 0              & 4,00          \\ \hline
2         & 1             & 0              & 5,50          \\ \hline
3         & 1,5           & 5,75           & 5,75          \\ \hline
4         & 2             & 6,50           & 4,25          \\ \hline
\end{tabular}
\end{center}
\end{table}


\subsection*{Exercice 4}
Trouvez le taux de rendement exigés \textbf{y} qui correspond à une échéance de 0,5 an.

\begin{enumerate}[a)]
\item 0.05757
\item 0.055
\item 0.04
\item 0.0425
\end{enumerate}

\subsection*{Exercice 5}
Trouvez le taux de rendement exigés \textbf{y} qui correspond à une échéance de 1 an.

\begin{enumerate}[a)]
\item 0.0575
\item 0.055
\item 0.04
\item 0.0425
\end{enumerate}

\subsection*{Exercice 6}
Trouvez le taux de rendement exigés \textbf{y} qui correspond à une échéance de 1.5 ans.

\begin{enumerate}[a)]
\item 0.0575
\item 0.055
\item 0.04
\item 0.0425
\end{enumerate}

\subsection*{Exercice 7}
Trouvez le taux de rendement exigés \textbf{y} qui correspond à une échéance de 2 ans.

\begin{enumerate}[a)]
\item 0.0575
\item 0.055
\item 0.04
\item 0.0425
\end{enumerate}

\section*{Section 3}
Considérer les informations suivantes sur des obligations gouvernementales: 

\begin{table}[H]
\begin{center}
\begin{tabular}{|l|l|l|l|}
\hline
Période t & Échéance (an) & TC (\%, m = 2) & y (\%, m = 2) \\ \hline
1         & 0,5           & 0              & 2,00          \\ \hline
2         & 1             & 0              & 4,00          \\ \hline
3         & 1,5           & 0              & 4,00          \\ \hline
4         & 2             & 2,00           & 2,00          \\ \hline
\end{tabular}
\end{center}
\end{table}

Calculer les taux de rendement au comptant (en version nominale annuelle capitalisée semestriellement) pour l'échéance:

\subsection*{Exercice 8}
Calculer le taux de rendement au comptant (en version nominale annuelle capitalisée semestriellement) pour l'échéance 0.5 an.

\begin{enumerate}[a)]
\item 1.975\%
\item 1.75\%
\item 2.00\%
\item 4.00\%
\end{enumerate}



\subsection*{Exercice 9}
Calculer le taux de rendement au comptant (en version nominale annuelle capitalisée semestriellement) pour l'échéance 1 an.

\begin{enumerate}[a)]
\item 1.975\%
\item 1.75\%
\item 2.00\%
\item 4.00\%
\end{enumerate}

\subsection*{Exercice 10}
Calculer le taux de rendement au comptant (en version nominale annuelle capitalisée semestriellement) pour l'échéance 1.5 ans.

\begin{enumerate}[a)]
\item 1.975\%
\item 1.75\%
\item 2.00\%
\item 4.00\%
\end{enumerate}


\subsection*{Exercice 11}
Calculer le taux de rendement au comptant (en version nominale annuelle capitalisée semestriellement) pour l'échéance 2 ans.

\begin{enumerate}[a)]
\item 1.975\%
\item 1.75\%
\item 2.00\%
\item 4.00\%
\end{enumerate}

\newpage

\section*{Section 4}
Considérer la structure à terme des taux d'intérêt suivante (les taux au comptant suivants):

\begin{table}[H]
\begin{center}
\begin{tabular}{|l|l|l|}
\hline
Période t & Échéance (an) & Taux au comptant effectif semestriel (\%) \\ \hline
1         & 0,5           & 1,00                                      \\ \hline
2         & 1             & 1,50                                      \\ \hline
3         & 1,5           & 2,00                                      \\ \hline
4         & 2             & 2,50                                      \\ \hline
\end{tabular}
\end{center}
\end{table}

\subsection*{Exercice 12}
Pour une obligation venant à échéance dans 2 ans (4 semestres) et ayant un taux de coupon semestriel de 4 \% et une valeur nominale de 100 \$,  trouver le prix théorique $P$

\begin{enumerate}[a)]
\item 105.00\$
\item 105.25\$
\item 105.63\$
\item 105.83\$
\end{enumerate}


\subsection*{Exercice 13}
Pour une obligation venant à échéance dans 2 ans (4 semestres) et ayant un taux de coupon semestriel de 4 \% et une valeur nominale de 100 \$,  trouver le taux de rendement à l'échéance y (en version nominale annuelle capitalisée semestriellement)

\begin{enumerate}[a)]
\item 4.954\%
\item 4.904\%
\item 4.854\%
\item 4.804\%
\end{enumerate}

\newpage

\section*{Section 5}
Considérer la structure à terme des taux d'intérêt suivante (les taux au comptant suivants): 

\begin{table}[H]
\begin{center}
\begin{tabular}{|l|l|l|}
\hline
Période t & Échéance (an) & Taux au comptant effectif semestriel (\%) \\ \hline
1         & 0,5           & 1,00                                      \\ \hline
2         & 1             & 1,00                                      \\ \hline
3         & 1,5           & 1,00                                      \\ \hline
4         & 2             & 1,50                                      \\ \hline
\end{tabular}
\end{center}
\end{table}

\subsection*{Exercice 14}
Trouver le taux de rendement à l'échéance au pair (en version nominale annuelle capitalisée semestriellement) pour l'échéance 0.5 an.

\begin{enumerate}[a)]
\item 1.00\%
\item 1.50\%
\item 2.00\%
\item 2.978\%
\end{enumerate}


\subsection*{Exercice 15}
Trouver le taux de rendement à l'échéance au pair (en version nominale annuelle capitalisée semestriellement) pour l'échéance 1 an.


\begin{enumerate}[a)]
\item 1.00\%
\item 1.50\%
\item 2.00\%
\item 2.978\%
\end{enumerate}

\subsection*{Exercice 16}
Trouver le taux de rendement à l'échéance au pair (en version nominale annuelle capitalisée semestriellement) pour l'échéance 1.5 ans.


\begin{enumerate}[a)]
\item 1.00\%
\item 1.50\%
\item 2.00\%
\item 2.978\%
\end{enumerate}



\subsection*{Exercice 17}
Trouver le taux de rendement à l'échéance au pair (en version nominale annuelle capitalisée semestriellement) pour l'échéance 2 ans.


\begin{enumerate}[a)]
\item 1.00\%
\item 1.50\%
\item 2.00\%
\item 2.978\%
\end{enumerate}


\end{document}
